\section*{Что вносит кластерное разложение C46}
В C46 энергия статической решётки и колебательная свободная энергия
аппроксимированы по 46 исходным декорациям, 20 кластерным функциям и
девяти объёмам при фиксированном \(c/a=4{,}6046/2{,}9446\). Сохранены
авторские маски обучения и фононной устойчивости. Из 409 устойчивых пар
«структура--объём» 346 входят в подгонку, 63 служат для контроля.
Для каждого спектра классическая квазигармоническая метка вычисляется как
\[
 F_{\rm cl}(T)=\frac{T}{T_0}F_{\rm cl}(T_0)-gk_{\rm B}T\ln(T/T_0),
 \qquad T_0=1000~\text{К},\quad g=\sum_\nu w_\nu.
\]
Здесь \(w_\nu\) --- веса мод в расчёте на атом. Для каждого спектра \(g\)
и \(F_{\rm cl}(T_0)\) восстановлены в расчёте при 1000~К из авторских
таблиц квантовой колебательной свободной энергии. Формула следует из
суммы \(k_{\rm B}T\ln(\hbar\omega_\nu/k_{\rm B}T)\) при фиксированном объёме.
При каждом \(T\) коэффициенты разложения заново подгоняются к меткам.
В конфигурационном Монте-Карло коэффициенты интерполируются по объёму,
который релаксирует при \(p=0\). На контрольных декорациях среднеквадратичная
ошибка аппроксимации колебательной свободной энергии по отдельным объёмам
лежит в диапазоне 0,25--1,71~мэВ/атом при 600~К и
0,53--3,63~мэВ/атом при 1273,15~К.

Для интегрирования по составу VCSGC в каждом окне к энергии добавляется
\(N[k_{\rm B}T\kappa(x-c_0)^2-\mu_0 x]\), \(\kappa=100\). Среднее
значение \(x\) даёт оценку производной свободной энергии для конечной ячейки,
\(g'(\langle x\rangle)\simeq\mu_0-2k_{\rm B}T\kappa(\langle x\rangle-c_0)\).
Интегрирование кусочно-кубического
интерполянта PCHIP, сохраняющего локальную монотонность, даёт \(g(x)\)
с точностью до константы.
Для пары фаз ищутся два внутренних контакта общей опорной касательной
с разнесением не менее 0,5~п.п. На 5120 узлах использованы 52 окна при
600~К и 51 при 1273,15~К, по 3000 проходов уравновешивания и 9000 сбора.

При 600~К классическая модель на 768 и 5120 узлах даёт соответственно
4,53/24,15 и 4,55/24,10~ат.\% Al. На большем размере условные интервалы
составляют [4,49; 4,61] и [24,08; 24,13]. Перевзвешивание тех же
конфигураций на квантовые колебательные коэффициенты даёт 4,60/24,09,
независимый авторский квантовый расчёт --- 4,59/24,09. Для \(\alpha\)
изменение при замене классических колебаний квантовыми составляет
\(+0{,}050\) п.п. с 95\%-м интервалом [0,026; 0,078], для \(\alpha_2\) ---
\(-0{,}010\) п.п. с интервалом [\(-0{,}024\); 0,004].

\textbf{Химические корреляции.} В приближении Брегга--Вильямса
с теми же классическими коэффициентами общая касательная даёт
2,47/24,28 при 600~К и 15,44/22,05 при 1273,15~К. После перехода к
квантовым коэффициентам получаются соответственно 2,49/24,27 и
15,49/22,07. На верхней изотерме среднее поле предсказывает ширину
6,61~п.п. Сопоставление с VCSGC показывает, насколько вывод о двухфазном
равновесии чувствителен к учёту корреляций между заполнениями узлов.

\subsection*{Следующий расчёт}
Для определения температурного интервала замыкания сольвуса целесообразно
добавить прямые изотермы от 1100 до 1250~К. На каждой нужны профили
порядка и состава для нескольких размеров при одинаковой форме образца.
Рост межфазного барьера и два сохраняющихся состава позволят проследить
линию первого рода. Вблизи исчезновения барьера нужны более крупные
ячейки и независимые затравки. Для различения локального порядка и
структурного перехода вблизи 27~ат.\% Al следует сохранять временные ряды
структурного фактора концентрационных волн D0$_{19}$ и сравнить по размерам
кумулянт Биндера и отношение корреляционной длины к размеру ячейки. Независимая
проверка разбавленных растворов и антиструктурных дефектов Ti$_3$Al по DFT
позволит оценить вклад самого MTP в низкотемпературное расхождение.

\medskip\noindent{\small\textbf{Архив расчёта.} В 34 заданиях UMT на аккаунте
\texttt{u3129} завершены 563 траектории прямого MTP. После прогрева сохранено
189\,150 кадров при 600~К и 286\,560 при 1273,15~К. Все архивы сверены
по SHA-256. Исходник отчёта включает численные данные обоих рисунков.}
